\section{IP Addressing, CIDR, and Subnetting Summary}

\subsection{IPv4 Address Structure}
An IPv4 address is a 32-bit value divided into two parts:
\begin{itemize}
    \item \textbf{Network portion} -- identifies the network
    \item \textbf{Host portion} -- identifies a device within that network
\end{itemize}

The prefix length (/N) specifies how many bits belong to the network. The remaining $(32 - N)$ bits represent the host.

\textbf{Key idea:} This split enables hierarchical routing, allowing routers to forward packets based on network prefixes instead of individual hosts. \\
\textit{(Prevents routing tables from becoming too large)} :contentReference[oaicite:0]{index=0}

\subsection{Classful Addressing (Deprecated)}
Originally, IP addresses were divided into classes:

\begin{itemize}
    \item Class A: /8
    \item Class B: /16
    \item Class C: /24
\end{itemize}

\textbf{Problems:}
\begin{itemize}
    \item Address space waste (allocations too large or too small)
    \item Routing table explosion (no aggregation)
\end{itemize}

This led to the adoption of CIDR.

\subsection{CIDR (Classless Inter-Domain Routing)}
CIDR allows flexible prefix lengths and efficient allocation.

\begin{itemize}
    \item Any prefix length is valid (e.g., /19, /27)
    \item Prefix explicitly written (e.g., 192.168.1.0/24)
    \item Enables route aggregation (supernetting)
\end{itemize}

\textbf{Subnet Mask:}
A subnet mask is a 32-bit value with $N$ ones followed by zeros.

\begin{center}
/24 $\rightarrow$ 255.255.255.0 \\
/26 $\rightarrow$ 255.255.255.192
\end{center}

\subsection{Network Calculations}
\begin{itemize}
    \item \textbf{Network Address:} IP $\AND$ subnet mask
    \item \textbf{Broadcast Address:} all host bits = 1
    \item \textbf{Usable Hosts:} $2^{(32-N)} - 2$
\end{itemize}

Example:
\begin{itemize}
    \item 192.168.1.0/24
    \item Network: 192.168.1.0
    \item Broadcast: 192.168.1.255
    \item Usable: 254 hosts
\end{itemize} :contentReference[oaicite:1]{index=1}

\subsection{Route Aggregation}
CIDR allows combining multiple networks into one larger prefix.

Example:
\begin{itemize}
    \item 192.168.0.0/24
    \item 192.168.1.0/24
    \item 192.168.2.0/24
    \item 192.168.3.0/24
\end{itemize}

Can be aggregated into:
\begin{center}
192.168.0.0/22
\end{center}

This reduces routing table size. :contentReference[oaicite:2]{index=2}

\subsection{Subnetting}
Subnetting divides a network into smaller sub-networks.

\textbf{Process:}
\begin{itemize}
    \item Borrow $k$ bits from host portion
    \item New prefix = $(N + k)$
    \item Number of subnets = $2^k$
\end{itemize}

Example: Split 192.168.10.0/24 into 4 subnets:
\begin{itemize}
    \item New prefix: /26
    \item Each subnet has 64 addresses (62 usable)
\end{itemize}

\textbf{Subnets:}
\begin{itemize}
    \item 192.168.10.0/26
    \item 192.168.10.64/26
    \item 192.168.10.128/26
    \item 192.168.10.192/26
\end{itemize} :contentReference[oaicite:3]{index=3}

\subsection{Subnet Membership}
To check if two hosts are in the same subnet:
\begin{itemize}
    \item Apply subnet mask to both
    \item Compare resulting network addresses
\end{itemize}

\subsection{VLSM (Variable Length Subnet Masking)}
Different subnets can have different sizes depending on need.

\begin{itemize}
    \item More efficient allocation
    \item Requires careful planning to avoid overlap
\end{itemize}

\subsection{Special IPv4 Address Ranges}
\begin{itemize}
    \item \textbf{Private (RFC 1918):}
    \begin{itemize}
        \item 10.0.0.0/8
        \item 172.16.0.0/12
        \item 192.168.0.0/16
    \end{itemize}
    \item \textbf{Loopback:} 127.0.0.0/8
    \item \textbf{Link-local:} 169.254.0.0/16
    \item \textbf{Multicast:} 224.0.0.0/4
\end{itemize} :contentReference[oaicite:4]{index=4}

\subsection{Address Allocation Hierarchy}
\begin{itemize}
    \item IANA
    \item Regional Internet Registries (RIRs)
    \item ISPs
    \item Organizations / End users
\end{itemize}

\subsection{IPv4 Exhaustion}
IPv4 has $2^{32} \approx 4.3$ billion addresses.

\begin{itemize}
    \item Exhausted globally by 2011--2021
    \item Addresses now bought/sold
\end{itemize}

\textbf{Solution:}
\begin{itemize}
    \item NAT -- allows many devices to share one IP
    \item IPv6 -- 128-bit addressing ($2^{128}$ addresses)
\end{itemize} :contentReference[oaicite:5]{index=5}

\subsection{Python \texttt{ipaddress} Module}
Python provides tools for subnet calculations:

\begin{verbatim}
import ipaddress
net = ipaddress.IPv4Network("192.168.10.0/24")
print(net.network_address)
print(net.broadcast_address)
print(net.hosts())
\end{verbatim}

Useful for:
\begin{itemize}
    \item Subnetting
    \item Address validation
    \item Network calculations
\end{itemize}